\subsection{庞加莱群的表现}

定理 1. --- 设 X 是一个紧的可解圈拓扑空间，x 是 X 的一点。庞加莱群 \(\pi_{1}(X,x)\) 是有限表现的。

x 在 X 中的弧连通分支是开的、闭的且可解圈的；这使我们可以假设空间 X 连通。由于空间 X 可解圈，X-空间 \(E = \lambda_{x}(X)\)——赋以终点映射——是一个非空且道路单连通的覆叠（IV, p. 342, 定理 1）。群 \(G = \text{Aut}_{\text{X}}(\text{E})\) 同构于 \(\pi_{1}(\text{X}, x)\)；因此只需证明群 G 是有限表现的。

由于 X 紧，X 的每一点 x 都有一个紧邻域 \(K_{x}\)，覆叠 E 在其上可平凡化（TG, I, p. 65, 推论）。由于 X 局部连通，每个 \(x \in X\) 都有一个开邻域 \(W_{x}\)（连通，且包含于 \(K_{x}\)）。设 F 是 X 的一个有限子集，使得诸 \(W_{x}\)（\(x \in F\)）覆盖 X。设 n 为 F 的基数。

我们用归纳法证明：对每个整数 \(k\)（\(1 \leqslant k \leqslant n\)），存在一个基数为 \(k\)、包含于 F 的子集 A，且对 \(x \in A\)，存在一个截面 \(s_x\)（\(p\) 在 \(K_x\) 上的截面），使得诸 \(s_x(W_x)\)（对 \(x \in A\)）之并是 E 的一个连通子集。断言对 \(k = 1\) 成立。假设它对一个整数 \(k\)（\(1 \leqslant k < n\)）成立，我们来证明它对 \(k + 1\) 成立。设 A 是 F 的基数为 \(k\) 的子集，且对每个 \(x \in A\)，设 \(s_x\) 是 E 在 \(K_x\) 上的一个截面，使 \(\bigcup_{x \in A} s_x(W_x)\) 连通。开集 \(\bigcup_{x \in A} W_x\) 与 \(\bigcup_{x \in F-A} W_x\) 非空且覆盖 X；由于 X 连通，它们的交非空。于是存在 \(x \in A\)、\(y \in F - A\) 及 \(z \in W_x \cap W_y\)。置 \(A' = A \cup \{y\}\)，并选取一个截面 \(s_y\)（E 在 \(K_y\) 上的截面），使 \(s_y(z) = s_x(z)\)。连通开集 \(\bigcup_{p \in A} s_p(W_p)\) 与 \(s_y(W_y)\) 非空且有公共点；因此其并连通。这就对 \(k + 1\) 证明了断言。由归纳法，它于是对每个整数 \(k \in \{1, \ldots, n\}\) 成立。

把上述结果应用于 k = n，于是对每个 \(x \in F\)，存在一个截面 \(s_x\)（E 在 \(K_x\) 上的截面），使得 \(U = \bigcup_{x \in F} s_x(W_x)\) 是 E 的一个开连通子集。我们有 \(p(U) = X\)。由于群 G 在覆叠 E 的每条纤维中传递地作用，我们有 GU = E。

U 的闭包包含于 \(\bigcup_{x\in\mathrm{F}}s_x(\mathrm{K}_x)\)，因而是紧的。G 在 E 上的作用是逆紧的（I, p. 96, 推论 1），故由偶 \((g,x)\in\mathrm{G}\times\mathrm{E}\)（满足 \(x\in\overline{\mathrm{U}}\) 且 \(gx\in\overline{\mathrm{U}}\)）所成之集是 \(\mathrm{G}\times\mathrm{E}\) 的一个紧子集（TG, I, p. 77, 命题 6）。由此得出，由 \(g\in\mathrm{G}\)（满足 \(\overline{\mathrm{U}}\cap g\overline{\mathrm{U}}\neq\varnothing\)）所成之集是紧的。由于它离散，它有限。更有，由 \(g\in\mathrm{G}\)（满足 \(\mathrm{U}\cap g\mathrm{U}\neq\varnothing\)）所成之集有限。定理于是由 I, p. 136 的命题 10 得出。

